\documentclass[10pt,a4paper]{amsart}
\usepackage[margin=19mm]{geometry}
\usepackage{amsmath,amssymb,mathtools}
\usepackage[scr=rsfs,cal=euler]{mathalfa}
\usepackage{xfrac}
\usepackage[unicode,hidelinks,bookmarks=false]{hyperref}
\newcommand{\ie}{i.e.\ }
\newcommand{\cf}{cf.\ }
\newcommand{\resp}{respectively\ }
\newcommand{\Subsection}{\subsection}
\providecommand{\nobreakdash}{\nobreak\hbox{-}}
\setlength{\emergencystretch}{2em}
\multlinegap0pt
\usepackage{bm}
\usepackage{bbm}
\usepackage{dsfont}
\usepackage{xspace}
\usepackage{cancel}
\usepackage{mathtools}
\usepackage{amsthm}
\makeatletter
\@ifundefined{theorem}{\newtheorem{theorem}{Theorem}}{}
\makeatother
\DeclareMathOperator{\Spec}{Spec}
\newcommand{\N}{\mathbb{N}}
\newcommand{\cA}{\mathcal{A}}
\newcommand{\gp}{\mathrm{gp}}
\title[The logarithmic $h$- and $v$-topologies]{The logarithmic $h$- and $v$-topologies}
\author{Nikolai Opdan, Doosung Park, Paul Arne {\O}stv{\ae}r}
\date{Source: arXiv:2608.06882v1 (CC BY 4.0)}
\begin{document}
\maketitle

\noindent\emph{Source and licence.} The logarithmic $h$- and $v$-topologies. Version arXiv:2608.06882v1 (CC BY 4.0), distributed under the Creative Commons Attribution 4.0 International licence, as recorded in the arXiv licence metadata for this exact version and shown on the abstract page https://arxiv.org/abs/2608.06882v1. Full rights notice: licenses/arxiv-2608.06882-CC-BY-4.0.txt.\par\smallskip

\noindent\textit{Editorial scope.} Counted unit: the Theorem whose source label is \texttt{thm:logvaluationringsliftspecializations} in the article (source0.tex lines 612--621, source label \texttt{thm:logvaluationringsliftspecializations}), delivered together with the article's own complete original proof of it (source0.tex lines 622--695, one \texttt{proof} environment of 74 lines). No mathematics is written, repaired or re-derived: statement and proof are the article's own text line for line. Required context retained uncounted: none. Every definition, display and label of those lines is retained unchanged. Omitted from this package and not redistributed: 1--611, 696--2211. Nothing inside a retained excerpt is omitted or altered: every retained body is delivered exactly as the article's lines are written, so no omission receipt of any kind accompanies this package. Citations: the retained excerpts carry no citation command at all, so no bibliography entry of the article is transcribed and no reference is redistributed. Reference adaptation: none was needed and none was made; every reference inside the delivered unit resolves inside it, so no unresolved reference or citation is produced. Printed slips kept literal: no printed word, symbol, hypothesis or number of the counted statement or of its proof was altered. Layout and class: the article's own class is \texttt{amsart} with packages including xcolor, hyperref, amsmath, amsthm, amssymb, mathrsfs, mathtools, tikz-cd, todonotes, comment, thmtools, enumitem, varioref, cleveref; the wrapper is the lane's pinned \texttt{amsart} editorial wrapper at 10pt on A4, which restates the theorem-like environments the retained text needs and adds the article's own macro definitions that the retained text needs (\texttt{\textbackslash N}, \texttt{\textbackslash Spec}, \texttt{\textbackslash cA}, \texttt{\textbackslash gp}), with extra wrapper packages (bm, bbm, dsfont, xspace, cancel, mathtools). The delivered numbering is the unit's own. Assets: no retained span contains a figure environment, a graphics-inclusion command, a PDF-page inclusion, a table or a TikZ picture; no image file is copied into this package, the package declares no asset dependency and no image file is redistributed with it. Licence: version arXiv:2608.06882v1 of this article is distributed under the Creative Commons Attribution 4.0 International licence, as recorded in the arXiv licence metadata for this exact version and shown on the abstract page https://arxiv.org/abs/2608.06882v1. A full rights notice is delivered at licenses/arxiv-2608.06882-CC-BY-4.0.txt. Characters: every retained excerpt is pure ASCII, so the delivered engine, XeLaTeX with its default Latin Modern text font, reports no missing character.
\begin{theorem}\label{thm:logvaluationringsliftspecializations}
Let $X$ be a quasi-coherent integral log scheme, 
i.e., $X$ admits integral charts Zariski locally.
Suppose se that $x' \rightsquigarrow x$ is a specialization
of points in $X$. 
Then there exists a log valuation ring $(V,P)$, 
with maximal ideal $\mathfrak{m} \subset V$, 
and a morphism $\Spec{(V,P)} \to X$ such that the specialization
$\mathfrak{m} \rightsquigarrow (0)$ lifts $x' \rightsquigarrow x$.
\end{theorem}

\begin{proof}
We may assume $X$ is the spectrum of a log ring $(A,M),$ 
where $A$ is an integral local ring and $M$ is a monoid, 
and where the generic (resp.\ unique closed) point of $X$ is $x'$ (resp.\ $x$).
Let $K$ be the fraction field of $A$.
A valuation ring $V$ exists with fraction field $K$ dominating $A$.
We need to show that the homomorphism of monoids $\theta\colon M\to V$ 
factors through a valuative monoid $P$.
It will be convenient to use additive instead of multiplicative notation for the 
monoid operation in a ring.
To avoid confusion,
we write $Q$ for the monoid $V$ in additive notation
and let $*$ be the element of $Q$ corresponding to $0\in V$.
Observe that $*+q=*$ for every $q\in Q$.

Let $\cA$ be the set of extensions $N\to Q$ of $\theta$ such that $N$ is an integral monoid,
$N$ dominates $M$,
and $M^\gp\to N^\gp$ is an isomorphism.
If $\{N_\alpha\to Q\}$ is a chain in $\cA$,
then $\bigcup_\alpha N_\alpha \to Q$ is an element of $\cA$.
Hence, 
by the Hausdorff maximality principle, 
a maximal element $\eta\colon P\to Q$ of $\cA$ exists.

We claim that $P$ is valuative.
If not, an element $x\in P^\gp$ exists such as $x,-x\notin P$.
Consider the face $F:=\eta^{-1}(Q-\{*\})$ of $P$,
and let $\alpha\colon F\to Q-\{*\}$ be the restriction of $\eta$.
First, assume $x\in F^\gp$.
Since $Q-\{*\}$ is valuative,
at least one of $\alpha^\gp(x)$ and $\alpha^\gp(-x)$ is in $Q-\{*\}$,
and we may assume $q:=\alpha^\gp(x)\in Q-\{*\}$ without loss of generality.
Let $P'$ be the submonoid of $P^{\gp}$ generated by $P$ and $x$.
We claim that $\eta'\colon P'\to Q$ given by $\eta'(p+nx):=\eta(p)+n*$ for $p\in P$ and $n\in \N$ is well-defined.
Assume that we have $p+nx=p'+n'x$ with $p'\in P$ and $n'\in \N$.
If $n,n'>0$,
then $\eta'(p+nx)=*=\eta'(p'+n'x)$.
In all other cases, we may assume $n'=0$ without loss of generality.
If $p\notin F$, then $p\notin F^\gp$,
so $p'\notin F$ since $x\in F^\gp$.
Hence $\eta'(p+nx)=\eta'(p')=*$.
If $p\in F$,
then $p'\in F$ since $x\in F^\gp$.
Thus, 
$\eta'(p+nx)=\eta'(p')=\alpha^\gp(p+nx)$, 
so that $\eta'$ is well-defined.
Hence $\eta'$ is an extension of $\eta$,
which contradicts the maximality of $\eta$.

Next, we assume $x\notin F^\gp$.
If $p+nx,p'-n'x\in F$ for some $p,p'\in P$ and $n,n'\in \N^+$,
then $n'p+np'\in F$,
so $p\in F$.
This implies $n(x+p)\in F$.
Since $F$ is saturated,
we have $x+p\in F$ and hence $x\in F^\gp$,
which is a contradiction.
Hence we may assume that $p+nx\notin F$ for every $p\in P$ and $n\in \N^+$.
Let $P'$ be the submonoid of $P^\gp$ generated by $P$ and $x$.
Next we claim that $\eta'\colon P'\to Q$ given by $\eta'(p+nx):=\eta(p)+n\ast$ is well-defined.
Assume that $p+nx=p'+n'x$, where $p'\in P$ and $n\in \N$.
If $n,n'>0$,
then $\eta'(p+nx)=*=\eta'(p'+n'x)$.
In all other cases, we assume $n'=0$ without loss of generality.
If $n=0$,
then we have $\eta'(p+nx)=\eta'(p')$.
If $n>0$,
then $\eta'(p+nx)=*$.
On the other hand,
the assumption on $x$ implies $p'\notin F$.
Thus, $\eta(p')=*$, so that $\eta'$ is well-defined.
Hence $\eta'$ extends $\eta$,
contradicting the maximality of $\eta$.
\end{proof}

\end{document}
